\documentclass[twocolumn,fleqn]{jsarticle}
\usepackage{ascmac,amsmath,amsthm,amssymb}
\setlength{\columnseprule}{0.4pt}
\setlength{\mathindent}{0mm}
\pagestyle{empty}
\begin{document}
\begin{itembox}[l]{項別微分}
関数列$\{f_n\}が区間[a,b]$上で以下の条件を満たすとき、項別微分が可能となる。
\begin{enumerate}
\item $\{f_n\}は[a,b]上のC^1級関数列$
\item $\displaystyle \sum_{n=1}^{\infty} f_n(x) = S(x)$ (各点収束)
\item $\displaystyle \sum_{n=1}^{\infty} f'_n(x)はT(x)$に一様収束
\end{enumerate}
\end{itembox}
\begin{proof}
$S_n(x):=\displaystyle \sum_{k=1}^{n} f_k(x)$とおく。このとき、

$f_nがC^1級であることよりS'_n(x)は[a,b]上で連続なので、[a,b]で$Riemann積分可能となる。

よって、$S_n(x)=S_n(a)+\displaystyle \int_a^x S'_n(t) dt$──(※)と表せる。

また、$S'_n(x)=\displaystyle \sum_{k=1}^{n} f'_k(x)はn \to {\infty}でT(x)$に一様収束するので、
$\displaystyle \int_a^x S'_n(t) dt はn \to {\infty}で\displaystyle \int_a^x T(t) dt$に一様収束する(示せる)

よって、(※)の両辺を$n \to {\infty}$とすると、

\bigskip

$S(x) = S(a) + \displaystyle \int_a^x T(t) dt$

\bigskip

両辺を$x$で微分すると、

$S'(x) = T(x)$, 即ち $(\displaystyle \sum_{n=1}^{\infty} f_n(x))'=$
$\displaystyle \sum_{n=1}^{\infty} f'_n(x)$
\end{proof}

\begin{itembox}[l]{項別積分}
$[a,b]上の連続関数列\{f_n\}に対して、S_n(x):=\displaystyle \sum_{k=1}^{n} f_k(x)が$
$n \to {\infty}でS:=\displaystyle \sum_{k=1}^{\infty} f_k(x)$ に一様収束するとき、項別積分が可能となる。
\end{itembox}

\begin{proof}
$f_n:連続 \ より、[a,b]$上でRiemann積分可能。このとき、
$\displaystyle \sum_{k=1}^{n} \displaystyle \int_a^b f_k(x) dx=\displaystyle \int_a^b S_n(x) dx$
において、$S_nはn \to {\infty}でSに一様収束することから\displaystyle \int_a^b S_n(x) dx$は
$n \to {\infty}で\displaystyle \int_a^b S(x) dx$に一様収束する(示せる）ので、両辺を
$n \to {\infty}$とすると、

$\displaystyle \sum_{k=1}^{\infty} \displaystyle \int_a^b f_k(x) dx$=
$\displaystyle \int_a^b S(x) dx$
\end{proof}

\end{document}
